%------------------------------------------------------------------------------%
% \chapter{Matrizes}
%------------------------------------------------------------------------------%

%------------------------------------------------------------------------------%
\begin{exercicios}{Matrizes}
    
    \begin{exercicio}
        Indicar explicitamente os elementos da matriz $A=(a_{ij})_{2\times 3}$ tal que $a_{ij} = i - j$
        \begin{answers}
            \(\qquad A = \left[
                \begin{array}{rrr}
                	0 & -1 & -2 \\
                	1 &  0 & -1
                \end{array}
                \right]
            \)
        \end{answers}
    \end{exercicio}
    
    \begin{exercicio}
        Ache a matriz $A$ do tipo $2 \times 3$ definida por $a_{ij} = i j$ onde $i$ indica a linha e $j$, a coluna.
        \begin{answers}
            \( \qquad A = \left[
            \begin{array}{rrr}
            	1 & 2 & 3 \\
            	2 & 4 & 6
            \end{array}
            \right]
            \)
        \end{answers}
    \end{exercicio}
    
    \begin{exercicio}
        Construir as seguintes matrizes
        \begin{alphalist}
            \item $A=(a_{ij})_{3\times 3}$ tal que 
                  \( a_{ij} =\begin{cases}
                                1, & \text{se } i=j \\
                                0, & \text{se } i\neq j
                            \end{cases}  \)
            \item $B=(b_{ij})_{3\times 3}$ tal que 
                  \( b_{ij} =\begin{cases}
                                1, & \text{se } i+j=4 \\
                                0, & \text{se } i+j\neq 4
                            \end{cases}  \)
        \end{alphalist}
        \begin{answers}
            \begin{mathlist}[2]
                \item A = 
                    \left[
                    \begin{array}{rrr}
                    	1 & 0 & 0 \\
                    	0 & 1 & 0 \\
                    	0 & 0 & 1
                    \end{array}
                    \right]
                \item A = 
                    \left[
                    \begin{array}{rrr}
                    	0 & 0 & 1 \\
                    	0 & 1 & 0 \\
                    	1 & 0 & 0
                    \end{array}
                    \right]
            \end{mathlist}
        \end{answers}
    \end{exercicio}
    
    \begin{exercicio}
        Construa a matriz $A = (a_{ij})_{2 \times 3}$ definida por
        \[a_{ij} = \begin{cases}
            3i + j,  & \text{se } i < j \\
            7,       & \text{se } i = j \\
            i^2 + j, & \text{se } i > j
        \end{cases}\]
    \end{exercicio}
    
    \begin{exercicio}
        Dada as matrizes
        \[A = \left[\begin{array}{cc}
            x + y &   1   \\
             -5   & x - y
        \end{array} \right]
        \qquad
        B = \left[\begin{array}{rr}
             3 &  1 \\
            -5 & -1
        \end{array} \right]\]
        calcule $x$ e $y$ de modo que $A=B$.
    \end{exercicio}
        
\end{exercicios} % Matrizes

%------------------------------------------------------------------------------%
\begin{exercicios}{Operações com Matrizes}
        
    \begin{exercicio}
        Datas as matrizes 
        \[ A = \left[\begin{array}{cc}
            5 & 6 \\
            4 & 2
        \end{array} \right]
        \qquad
        B = \left[\begin{array}{rr}
            0 & -1 \\
            5 & 4
        \end{array} \right] \]
        calcule $A+B$, $A-B$, $2A+3B$, $AB$, $BA$, $A^t$ e $A^t+B$.
        \begin{answers}
            \begin{mathlist}[2]
                \item A+B   =   \left[
                \begin{array}{rr}
                    5 &  5 \\
                    9 &  6 \\
                \end{array}
                \right]
                \item A-B   =   \left[
                \begin{array}{rr}
                    5 &  7 \\
                    -1 & -2 \\
                \end{array}
                \right]
                \item 2A+3B =   \left[
                \begin{array}{rr}
                    10 &  9 \\
                    23 & 16 \\
                \end{array}
                \right]
                \item AB    =   \left[
                \begin{array}{rr}
                    30 & 19 \\
                    10 &  4 \\
                \end{array}
                \right]
                \item BA    =   \left[
                \begin{array}{rr}
                    -4 & -2 \\
                    41 & 38 \\
                \end{array}
                \right]
                \item A^t   =   \left[
                \begin{array}{rr}
                    5 &  4 \\
                    6 &  2 \\
                \end{array}
                \right]
                \item A^t+B =   \left[
                \begin{array}{rr}
                    5 &  3 \\
                    11 &  6 \\
                \end{array}
                \right]
            \end{mathlist}
        \end{answers}
    \end{exercicio}
    
    \begin{exercicio}
        Datas as matrizes
        \[ A = \left[\begin{array}{rrr}
            1 & 5 & 7 \\
            3 & 9 & 11
        \end{array} \right]
        \]\[
        B = \left[\begin{array}{rrr}
            2 &  4 &  6 \\
            8 & 10 & 12
        \end{array} \right]
        \]\[
        C = \left[\begin{array}{rrr}
            0 & -1 & -5 \\
            1 &  4 &  7
        \end{array} \right] \]
        calcule, se possível, $A+B+C$, $A-B+C$, $A-B-C$, $-A+B-C$, $AB$, $BC$, $AB^t$, $CB^t$, $A^tB$ e $A^tB^t$.
        \begin{answers}
            \begin{mathlist}[2]
            \item A+B+C    =   \left[
            \begin{array}{rrr}
                3 &  8 &  8 \\
                12 & 23 & 30 \\
            \end{array}
            \right]
            \item A-B+C    =   \left[
            \begin{array}{rrr}
                -1 &  0 & -4 \\
                -4 &  3 &  6 \\
            \end{array}
            \right]
            \item A-B-C    =   \left[
            \begin{array}{rrr}
                -1 &  2 &  6 \\
                -6 & -5 & -8 \\
            \end{array}
            \right]
            \item -A+B-C   =   \left[
            \begin{array}{rrr}
                1 &  0 &  4 \\
                4 & -3 & -6 \\
            \end{array}
            \right]
            \item AB \) Não pode ser calculada
            \item BC \) Não pode ser calculada
            \item AB^t     =   \left[
            \begin{array}{rr}
                64 & 142 \\
                108 & 246 \\
            \end{array}
            \right]
            \item CB^t     =   \left[
            \begin{array}{rr}
                -34 & -70 \\
                60 & 132 \\
            \end{array}
            \right]
            \item A^tB     =   \left[
            \begin{array}{rrr}
                26 & 34 & 42 \\
                82 & 110 & 138 \\
                102 & 138 & 174 \\
            \end{array}
            \right]
            \item A^tB^t \) Não pode ser calculada
            \end{mathlist}
        \end{answers}
    \end{exercicio}

    \begin{exercicio}
        Calcule os seguintes produtos matriciais
        \begin{mathlist}
            \item \left[\begin{array}{rrr}
                  1 & 2 & 3 \\
                  4 & 5 & 6
                \end{array} \right] 
                \left[\begin{array}{r}
                    7 \\ 8 \\ 9
                \end{array}\right]
            \item \left[\begin{array}{rr}
                   0 & 1 \\
                   1 & 0
                \end{array} \right] 
                \left[\begin{array}{rr}
                   4 & 7 \\ 
                   2 & 3
                \end{array}\right]
            \item \left[\begin{array}{rrr}
                    1 & 5 & 2 \\
                   -1 & 4 & 7
                \end{array} \right] 
                \left[\begin{array}{rr}
                    1 & -1 \\ 
                    2 &  3 \\
                   -3 &  0
                \end{array}\right]
            \item \left[\begin{array}{r}
                    1 \\ 2 \\ 3 
                \end{array} \right] 
                \left[\begin{array}{rrrr}
                    3 & 1 & 1 & 2
                \end{array}\right]
            \item \left[\begin{array}{rrrr}
                   1 & -1 & 5 & 0 \\
                   2 &  3 & 7 & 1 
                \end{array} \right] 
                \left[\begin{array}{rr}
                   1 & 1 \\ 
                   2 & 1 \\
                   3 & 1 \\
                   1 & 1
                \end{array}\right]
            \item \left[\begin{array}{rrr}
                   0 & 1 & 1 \\
                   2 & 2 & 0 \\
                   0 & 3 & 4 
                \end{array} \right] 
                \left[\begin{array}{rrr}
                   1 & 4 & 7 \\ 
                   0 & 0 & 1 \\
                   1 & 2 & 0
                \end{array}\right]
        \end{mathlist}
        \begin{answers}
            \begin{mathlist}[2]
                \item   \left[
                \begin{array}{r}
                    50 \\
                    122 \\
                \end{array}
                \right]
                \item   \left[
                \begin{array}{rr}
                    2 &  3 \\
                    4 &  7 \\
                \end{array}
                \right]
                \item   \left[
                \begin{array}{rr}
                    5 & 14 \\
                    -14 & 13 \\
                \end{array}
                \right]
                \item   \left[
                \begin{array}{rrrr}
                    3 &  1 &  1 &  2 \\
                    6 &  2 &  2 &  4 \\
                    9 &  3 &  3 &  6 \\
                \end{array}
                \right]
                \item   \left[
                \begin{array}{rr}
                    14 &  5 \\
                    30 & 13 \\
                \end{array}
                \right]
                \item   \left[
                \begin{array}{rrr}
                    1 &  2 &  1 \\
                    2 &  8 & 16 \\
                    4 &  8 &  3 \\
                \end{array}
                \right]
            \end{mathlist}
        \end{answers}
    \end{exercicio}

    \begin{exercicio}
        Sendo \( A=\left[\begin{array}{rr}
                                    1 & 1 \\
                                    0 & 1
                                 \end{array}\right] \)
        calcule $A^2$, $A^3$, $A^4$ e $A^n$.
        \begin{answers}
            \begin{mathlist}[2]
                \item A^2 =   \left[
                \begin{array}{rr}
                    1 &  2 \\
                    0 &  1 \\
                \end{array}
                \right]
                \item A^3 =   \left[
                \begin{array}{rr}
                    1 &  3 \\
                    0 &  1 \\
                \end{array}
                \right]
                \item A^4 =   \left[
                \begin{array}{rr}
                    1 &  4 \\
                    0 &  1 \\
                \end{array}
                \right]
                \item A^n =   \left[
                \begin{array}{rr}
                    1 &  n \\
                    0 &  1 \\
                \end{array}
                \right]
            \end{mathlist}
        \end{answers}
    \end{exercicio}
        
    \begin{exercicio}
        Calcule $AB$, $BA$, $A^2$ e $B^2$ sabendo que
        \[ A=\left[\begin{array}{rr}
                     2 &  1 \\
                    -4 & -2
                    \end{array}\right]
           \qquad
           B=\left[\begin{array}{rr}
                    2 & 1 \\
                    1 & 0
                   \end{array}\right]
        \]
        \begin{answers}
            \begin{mathlist}[2]
                \item AB  =   \left[
                \begin{array}{rr}
                    5 &  2 \\
                    -10 & -4 \\
                \end{array}
                \right]
                \item BA  =   \left[
                \begin{array}{rr}
                    0 &  0 \\
                    2 &  1 \\
                \end{array}
                \right]
                \item A^2 =   \left[
                \begin{array}{rr}
                    0 &  0 \\
                    0 &  0 \\
                \end{array}
                \right]
                \item B^2 =   \left[
                \begin{array}{rr}
                    5 &  2 \\
                    2 &  1 \\
                \end{array}
                \right]
            \end{mathlist}
        \end{answers}
    \end{exercicio}
        
    \begin{exercicio}
        Calcular o produto $ABC$ sendo dados
        \[   A=\left[\begin{array}{rr}
                 1 & 2 \\
                 5 & 1
                \end{array}\right]
            \hfill
            B=\left[\begin{array}{rrr}
                1 & 1 & 1\\
                3 & 2 & 1
                \end{array}\right]
            \hfill
            C=\left[\begin{array}{rr}
                3 &  1 \\
                1 &  0 \\
                2 & -1
                \end{array}\right]
        \]
        \begin{answers}
            \[ ABC =   \left[
                \begin{array}{rr}
                	32 & 4 \\
                	43 & 2
                \end{array}
                \right]
            \]
        \end{answers}
    \end{exercicio}
    
   \begin{exercicio}
       Considerando as matrizes
       \[ A = \left[
       \begin{array}{rrr}
       1 & 2 & -3 \\
       3 & 4 &  0
       \end{array} 
       \right]
       \qquad
       B = \left[
       \begin{array}{rrr}
       -2 & 1 &  5 \\
       0 & 3 & -4
       \end{array} 
       \right]
       \]
       Calcule
       \begin{mathlist}[2]
           \item C = A + B
           \item D = A - B
           \item E = 2A + 3B
           \item F = A^t + B^t
           \item G = A B^t
           \item H = A^t B
        \end{mathlist}
        \begin{answers}
            \begin{mathlist}[2]
                \item C =   \left[
                \begin{array}{rrr}
                    -1 &  3 &  2 \\
                    3 &  7 & -4 \\
                \end{array}
                \right]
                \item D =   \left[
                \begin{array}{rrr}
                    3 &  1 & -8 \\
                    3 &  1 &  4 \\
                \end{array}
                \right]
                \item E =   \left[
                \begin{array}{rrr}
                    -4 &  7 &  9 \\
                    6 & 17 & -12 \\
                \end{array}
                \right]
                \item F =   \left[
                \begin{array}{rr}
                    -1 &  3 \\
                    3 &  7 \\
                    2 & -4 \\
                \end{array}
                \right]
                \item G =   \left[
                \begin{array}{rr}
                    -15 & 18 \\
                    -2 & 12 \\
                \end{array}
                \right]
                \item H =   \left[
                \begin{array}{rrr}
                    -2 & 10 & -7 \\
                    -4 & 14 & -6 \\
                    6 & -3 & -15 \\
                \end{array}
                \right]
            \end{mathlist}
        \end{answers}
    \end{exercicio}
    
    \begin{exercicio}
        Considerando as matrizes
        \[ A = \left[ 
            \begin{array}{rrr}
               	1 & 2 & -3 \\
               	3 & 4 &  0
            \end{array} 
            \right]
            \qquad
            B = \left[
            \begin{array}{rrr}
               	-2 &  1 & 0 \\
              	 0 &  3 & 0 \\
              	 5 & -4 & 0
            \end{array} 
            \right]
        \]
        Calcule
        \begin{mathlist}[2]
            \item C = AB
            \item D = A(3B)
            \item E = B^2
            \item F = B^3
            \item G = B B^t
            \item H = A^t A
            \item J = A A^t
        \end{mathlist}
        \begin{answers}
            \begin{mathlist}[2]
                \item C =   \left[
                \begin{array}{rrr}
                    -17 & 19 &  0 \\
                    -6 & 15 &  0 \\
                \end{array}
                \right]
                \item D =   \left[
                \begin{array}{rrr}
                    -51 & 57 &  0 \\
                    -18 & 45 &  0 \\
                \end{array}
                \right]
                \item E =   \left[
                \begin{array}{rrr}
                    4 &  1 &  0 \\
                    0 &  9 &  0 \\
                    -10 & -7 &  0 \\
                \end{array}
                \right]
                \item F =   \left[
                \begin{array}{rrr}
                    -8 &  7 &  0 \\
                    0 & 27 &  0 \\
                    20 & -31 &  0 \\
                \end{array}
                \right]
                \item G =   \left[
                \begin{array}{rrr}
                    5 &  3 & -14 \\
                    3 &  9 & -12 \\
                    -14 & -12 & 41 \\
                \end{array}
                \right]
                \item H =   \left[
                \begin{array}{rrr}
                    10 & 14 & -3 \\
                    14 & 20 & -6 \\
                    -3 & -6 &  9 \\
                \end{array}
                \right]
                \item J =   \left[
                \begin{array}{rr}
                    14 & 11 \\
                    11 & 25 \\
                \end{array}
                \right]
            \end{mathlist}
        \end{answers}
    \end{exercicio}
    
    \begin{exercicio}
        Considere as seguintes matrizes
        \[
        A = \left[\begin{array}{cc}
           	2 & 0 \\
           	6 & 7
        \end{array}\right]
        \qquad
        B = \left[\begin{array}{rr}
           	0 &  4 \\
           	2 & -8
        \end{array}\right]
        \]\[
        C = \left[\begin{array}{rrr}
           	-6 &  9 & -7 \\
          	 7 & -3 & -2
        \end{array}\right]
        \]\[ 
        D = \left[\begin{array}{rrr}
           	-6 & 4 & 0 \\
          	 1 & 1 & 4 \\
           	-6 & 0 & 6
        \end{array}\right] 
        \qquad
        E = \left[\begin{array}{rrr}
          	 6 & 9 & -9 \\
           	-1 & 0 & -4 \\
           	-6 & 0 & -1
        \end{array}\right] 
        \]
        Se for possível calcule
        \begin{mathlist}[2]
            \item AB-BA
            \item 2C-D
            \item (2D^t-3E^t)^t
            \item D^2-DE
        \end{mathlist}
        \begin{answers}
            \begin{mathlist}[2]
                \item AB-BA         =   \left[
                \begin{array}{rr}
                    -24 & -20 \\
                    58 & 24 \\
                \end{array}
                \right]
                \item 2C-D \) Não pode ser calculada
                \item (2D^t-3E^t)^t =   \left[
                \begin{array}{rrr}
                    -30 & -19 & 27 \\
                    5 &  2 & 20 \\
                    6 &  0 & 15 \\
                \end{array}
                \right]
                \item D^2-DE        =   \left[
                \begin{array}{rrr}
                    80 & 34 & -22 \\
                    -10 & -4 & 45 \\
                    72 & 30 & -12 \\
                \end{array}
                \right]
            \end{mathlist}
        \end{answers}
    \end{exercicio}    

            
    \begin{exercicio}
        Sendo 
        \[ A=\left[\begin{array}{rr}
        	1 & 9 \\
        	2 & 1
        \end{array}\right]
        \qquad
        B=\left[\begin{array}{rr}
        	0 & -8 \\
        	3 & -1
        \end{array}\right]\]
        calcule
        \begin{mathlist}[2]
            \item (A+B)^2
            \item (A+B)(A-B)
            \item A^2 -2I_2A +I_2^2
            \item A^3 - I_2^3
        \end{mathlist}
        \begin{answers}
            \begin{mathlist}
                \item (A+B)^2           =   \left[
                \begin{array}{rr}
                    6 &  1 \\
                    5 &  5 \\
                \end{array}
                \right]
                \item (A+B)(A-B)        =   \left[
                \begin{array}{rr}
                    0 & 19 \\
                    5 & 85 \\
                \end{array}
                \right]
                \item A^2 -2I_2A +I_2^2 =   \left[
                \begin{array}{rr}
                    18 &  0 \\
                    0 & 18 \\
                \end{array}
                \right]
                \item A^3 - I_2^3       =   \left[
                \begin{array}{rr}
                    54 & 189 \\
                    42 & 54 \\
                \end{array}
                \right]
            \end{mathlist}
        \end{answers}
    \end{exercicio}
            
    \begin{exercicio}
        Sendo \( A=\left[\begin{array}{rr}
                                    1 & -1 \\
                                    0 &  2
                                 \end{array}\right] \),
        qual das matrizes comuta com $A$
        \begin{mathlist}[2]
            \item B = \left[\begin{array}{r}
                    2 \\ 3
                \end{array}\right]
            \item C = \left[\begin{array}{rrr}
                   1 & 3 & 2 \\
                   4 & 5 & 1
                \end{array}\right]
            \item D = \left[\begin{array}{rr}
                   0 & 0 \\
                   1 & 0
                \end{array}\right]
            \item E = \left[\begin{array}{rr}
                   5 & 2 \\
                   0 & 3
                \end{array}\right]
        \end{mathlist}
    \end{exercicio}
        
    \begin{exercicio}
        Calcule $x$, $y$, $z$ e $t$ sabendo que
        \[\left[\begin{array}{cc}
        x & 1 \\
        1 & 2
        \end{array}\right] +
        \left[\begin{array}{rr}
        2 &  y \\
        0 & -1
        \end{array}\right]  = 
        \left[\begin{array}{cc}
        3 & 2 \\
        z & t
        \end{array}\right]\]
    \end{exercicio}
    
    \begin{exercicio}
        Obtenha $X$ sabendo que 
        \[ X + \left[\begin{array}{r} 1 \\  4 \\  7 \end{array} \right]
        = \left[\begin{array}{r} 5 \\  7 \\  2 \end{array} \right]
        + \left[\begin{array}{r} 1 \\ -1 \\ -2 \end{array} \right]
        \]
    \end{exercicio}        
        
    \begin{exercicio}
        Determine $x$ e $y$ de modo que mas matrizes $A$ e $B$ comutem
        \[ A=\left[\begin{array}{rr}
                1 & 2 \\
                1 & 0
            \end{array}\right]
            \qquad
            B=\left[\begin{array}{rr}
                0 & 1 \\
                x & y
            \end{array}\right]
        \]
    \end{exercicio}
            
    \begin{exercicio}
        Obter todas as matrizes $B$ que comutem com 
          \[ A=\left[\begin{array}{rr}
                        1 & -1 \\
                        3 &  0
                    \end{array}\right]\]
    \end{exercicio}

    \begin{exercicio}
        Dada as matrizes
        \[A = \left[\begin{array}{r}
            -1 \\
             2
        \end{array} \right]
        \qquad
        B = \left[\begin{array}{c}
            3 \\
            0
        \end{array} \right]\] 
        \[C = \left[\begin{array}{r}
            -2 \\
             4
        \end{array} \right] 
        \qquad
        D = \left[\begin{array}{r}
             1 \\
            -1
        \end{array} \right]\]
        determine $X$ e $Y$, tais que
        \begin{mathlist}[1]
            \item X = A + 2B - C + D
            \item \frac{Y + A}{2} = \frac{C}{2} + 2B - 3D
        \end{mathlist}
    \end{exercicio}

    \begin{exercicio}
        Sendo
        \[A = \left[\begin{array}{cc}
            3 & 1 \\
            5 & 2
        \end{array} \right]
        \qquad
        B = \left[\begin{array}{cc}
            2 & 1 \\
            1 & 1
        \end{array} \right]\]
        \[I = \left[\begin{array}{cc}
            1 & 0 \\
            0 & 1
        \end{array} \right]    
        \qquad
        O = \left[\begin{array}{cc}
            0 & 0 \\
            0 & 0
        \end{array} \right]\]
        resolva o sistema, onde $X$ e $Y$ são matrizes quadradas de ordem 2
        \[\begin{cases}
            AX + Y = O \\
            BX + Y = I
        \end{cases}\]
    \end{exercicio}
    
    \begin{exercicio}
        Calcule $x$, data a igualdade
        \[\left[\begin{array}{cr}
            x & 1  \\
            3 & -1
        \end{array} \right]
        \left[\begin{array}{cc}
            2 & 0 \\
            1 & 1
        \end{array} \right]    =
        \left[\begin{array}{cr}
            0 & 1  \\
            5 & -1
        \end{array} \right]\]
    \end{exercicio}

    \begin{exercicio}
        Determine os valores de $x$, $y$ e $z$ na igualdade abaixo, envolvendo matrizes reais $2 \times 2$
        \[\left[\begin{array}{cc}
        0 & 0 \\
        x & 0
        \end{array} \right]    \,
        \left[\begin{array}{cc}
        0 & x \\
        0 & 0
        \end{array} \right] = 
        \left[\begin{array}{cc}
        x - y & 0 \\
        x   & z
        \end{array} \right] +
        \left[\begin{array}{cc}
        z - 4 & 0 \\
        y - z & 0
        \end{array} \right]\]
    \end{exercicio}
    
    \begin{exercicio}
        Calcule $x$ de modo que as matrizes
        \[A = \left[\begin{array}{rr}
        2 & -3 \\
        1 & 4
        \end{array} \right]    
        \qquad
        B = \left[\begin{array}{cc}
        4 & x \\
        0 & 4
        \end{array} \right]\]
        sejam comutativas.
    \end{exercicio}
    
    \begin{exercicio}
        Dada a matriz
        \[ A = \left[\begin{array}{cc}
        \sfrac{3}{4} & \sfrac{1}{2} \\[2mm]
        \sfrac{1}{4} & \sfrac{1}{2}
        \end{array} \right] \]
        determine uma matriz real da forma
        \[T = \left[\begin{array}{c}
        x   \\
        1 - x
        \end{array} \right]\]
        para a qual $AT = T$
    \end{exercicio}

    \begin{exercicio}
       Seja, para cada $x$ real, a matriz quadrada de ordem~2
       \[M (x) = \left[\begin{array}{rr}
            x & -x \\
           2x &  x
       \end{array} \right]\]
       Calcule o produto de $M(x)$ por $M(-x)$.
    \end{exercicio}
       
   \begin{exercicio}
       Se conhecemos apenas as matrizes $AB$ e $AC$ como podemos calcular as matrizes
       \begin{mathlist}[2]
           \item A(B+C)
           \item B^tA^t
           \item C^tA^t
           \item (ABA)C
       \end{mathlist}
    \end{exercicio}

    \begin{exercicio}
        Mostre que as matrizes da forma
        \[A = \left[ \begin{array}{ll}
            1 & y^{-1} \\
            y & 1
        \end{array} \right]\]
        em que $y$ é uma número real não nulo, verificam a equação \[X^2 = 2X\]
    \end{exercicio}
    
    \begin{exercicio}
        Verifique se as identidades são verdadeiras
        \begin{mathlist}
            \item (A+B)^2 = A^2 + 2AB + B^2
            \item (AB)C = C(AB)
            \item (ABC)^t = C^tB^tA^t
            \item (A+I)(A-I) = A^2 - I 
        \end{mathlist}
    \end{exercicio}

\end{exercicios} % Operações com Matrizes

%------------------------------------------------------------------------------%
\begin{exercicios}{Matriz Transposta}

    \begin{exercicio}
        Dada a matriz
        \[A = \left[\begin{array}{rrr}
        	 1 &  4 & 0 \\
        	-2 & -1 & 1 \\
        	 3 &  5 & 2
        \end{array} \right]    \]
        calcule
        \begin{mathlist}[2]
            \item X = A + A^t
            \item Y = A^t - A
        \end{mathlist}
        \begin{answers}
            \begin{mathlist}[2]
                \item X =   \left[
                \begin{array}{rrr}
                    2 &  2 &  3 \\
                    2 & -2 &  6 \\
                    3 &  6 &  4 \\
                \end{array}
                \right]
                \item Y =   \left[
                \begin{array}{rrr}
                	 0 & -6 & 3 \\
                	 6 &  0 & 4 \\
                	-3 & -4 & 0
                \end{array}
                \right]
            \end{mathlist}
        \end{answers}
    \end{exercicio}

    \begin{exercicio}
        Considere as matrizes
        \[A = \left[\begin{array}{rrr}
             a & b & 1 \\
            -1 & 1 & a
        \end{array} \right]
        \qquad
        B = \left[\begin{array}{rrr}
            1 & -1 & 0 \\
            0 &  1 & 0
        \end{array} \right]\]
        Determine $a$ e $b$, sabendo que
        \[A B^t = \left[\begin{array}{rr}
             3 & 4 \\
            -2 & 1
        \end{array} \right]\]
    \end{exercicio}

    \begin{exercicio}
        Sabendo que 
        \[ A = \left[ \begin{array}{r} x \\  4 \\ -2 \end{array} \right] \qquad
           B = \left[ \begin{array}{r} 2 \\ -3 \\  5 \end{array} \right] \]
        e que $A^tB=4$ determine o valor de $x$.
    \end{exercicio}

    \begin{exercicio}
        Encontre um valor de $x$ tal que $AB^t=0$, sendo que 
        \[ A = \left[ \begin{array}{ccc} x & 4 & -2 \end{array} \right]  \qquad
           B = \left[ \begin{array}{ccc} 2 & -3 & 5 \end{array} \right]
        \]
    \end{exercicio}
     
\end{exercicios} % Matriz Transposta
            
%------------------------------------------------------------------------------%
\begin{exercicios}{Matriz Inversa}
        
    \begin{exercicio}
        Mostre que $A^{-1}$ é a inversa de $A$
        \begin{mathlist}
            \item A = \left[\begin{array}{cc}
                     1 & 3 \\
                     2 & 7
                 \end{array} \right] 
                 \quad
                 A^{-1} = \left[\begin{array}{cc}
                     7  & -3 \\
                     -2 & 1
                 \end{array} \right]
            \item A = \left[\begin{array}{rrr}
                                        1 & 2 & 7 \\
                                        0 & 3 & 1 \\
                                        0 & 5 & 2
                                    \end{array} \right] \)
                  \( A^{-1} = \left[ \begin{array}{rrr}
                                            1 & 31 & -19 \\
                                            0 &  2 &  -1 \\
                                            0 & -5 &   3
                                        \end{array} \right] \)
        \end{mathlist}
    \end{exercicio}

    \begin{exercicio}
        Resolva a equação  matricial $AX = B$, dadas as matrizes
        \[A = \left[\begin{array}{rr}
             2 & -1 \\
            -3 &  2
        \end{array} \right]    
        \qquad
        B = \left[\begin{array}{rr}
             3 &  2 \\
            -6 & -5
        \end{array} \right]    \]
    \end{exercicio}

    \begin{exercicio}
        Determine $a$ real de modo que a matriz 
        \[A = \left[\begin{array}{cc}
            1 & 2 \\
            0 & a
        \end{array} \right]\]
        seja igual à sua inversa.
    \end{exercicio}
    
\end{exercicios} % Matriz Inversa

%------------------------------------------------------------------------------%
